% Literature review section (博士生1) — based on lit folder
% Standalone section: can be \input after \section{Literature Review} in main file.
% Citations refer to references in paper_draft/references.bib (lit folder entries added).

\subsection{Concepts and Measurement of Food Away from Home}

Food away from home (FAH), or out-of-home food consumption, refers to food prepared and consumed outside the household. Standard definitions in agricultural and consumer economics include meals and snacks purchased at restaurants, canteens, and other out-of-home establishments, and sometimes extend to institutional and business-provided meals \citep{USDA_ERS_FAH}. Measurement can be in \emph{expenditure} (spending on FAH) or in \emph{quantity} (physical consumption of food items consumed away from home). Expenditure-based measures are widely available from household budget and consumer expenditure surveys and have been used in a long line of work on FAH demand. Early US studies examined household consumption of food-away-from-home by total expenditure and by type of facility \citep{McCracken1987}, participation and expenditure patterns \citep{Byrne1996}, and the role of working wives and demographic factors \citep{Kinsey1983, Yen1993}. Similar approaches have been applied in other countries, for example panel data analysis of demand for food away from home in Spain \citep{Angulo2007}. Quantity-based measures require survey instruments that record both home and total (or out-of-home) consumption by food category, which are less common and harder to aggregate to the macro level.

For policy and food security analysis, a key quantity-based indicator is the ratio of total consumption to at-home consumption (or its inverse, the at-home share). In this paper, the \emph{out-of-home consumption coefficient} is defined as total consumption divided by at-home consumption, so that total food use can be inferred from at-home consumption data when the latter is the only dimension available in macro or regional statistics. This coefficient is distinct from FAH expenditure shares and elasticities commonly reported in the FAH literature and is the quantity-based input needed for food balance and demand projections.

\subsection{Empirical Evidence on China's Food Away from Home}

Evidence on China's FAH consumption comes mainly from household surveys and expenditure studies. \citet{Ma2006} showed that consumption of food away from home in urban China rises with income (``getting rich and eating out'') and documented strong income gradients using urban household data. \citet{Gould2006} assessed the current structure of food demand in urban China and provided estimates of demand elasticities that inform understanding of how food consumption responds to income and prices. Disaggregated evidence by type of facility and meal is available for Beijing: \citet{Bai2012} disaggregated household expenditures on food away from home in Beijing by type of food facility and type of meal, showing heterogeneity in where and how urban households consume food away from home. \citet{Liu2015} examined the role of household composition and income in food-away-from-home expenditure in urban China, reinforcing that demographics and income are key drivers. More recently, \citet{Zheng2019} predicted the changes in the structure of food demand in China, including the role of FAH in evolving diets. These studies establish that China's FAH is substantial and growing, that income and urbanization are primary drivers, and that household composition and demographics matter. They rely largely on expenditure or share measures from urban household surveys and the China Health and Nutrition Survey (CHNS); they do not produce nationally or provincially representative, commodity-level \emph{quantity} coefficients (total-to-at-home ratios) suitable for macro-level food balance and food security modeling.

\subsection{International Evidence and Methodological Approaches}

International work on FAH demand provides useful methodological benchmarks and confirms that the distinction between food at home and food away from home matters for demand and policy analysis. \citet{Balagtas2023} analyze consumer demand for food at home and food away from home and discuss responsiveness to policy. \citet{Brauw2020} address income variability and evolving diets in elasticity estimation for processed foods. Studies in other settings, such as consumption of processed food and food away from home in big cities, small towns, and rural areas \citep{Sauer2021}, show that settlement type and geography affect FAH patterns and that comparable quantity-oriented measures are scarce. Shocks such as the COVID-19 pandemic have been shown to affect food purchase behavior and food values \citep{Ellison2020}, underscoring the need for robust and updatable coefficient estimates. Together, this body of work supports the importance of FAH in modern food systems but also highlights that commodity-level, quantity-based out-of-home coefficients are rarely available for use in national or regional food balance and food security assessment.

\subsection{Data and Methodological Limitations}

Three limitations are central. First, \emph{micro--macro mismatch}: household surveys contain individual or household income and detailed consumption, whereas macro or provincial data typically provide only average income and aggregate or at-home consumption. Using micro elasticities or coefficients directly for macro prediction therefore requires a way to map macro covariates to the distribution of behavior observed in the micro data. Second, \emph{spatial and product coverage}: many surveys do not cover all provinces or all years, so that provincial-by-commodity out-of-home coefficients are missing for some region--year cells and must be inferred or interpolated. Third, \emph{measurement and model choice}: different surveys define ``at home'' and ``away from home'' differently, and the choice of functional form or model can affect estimated elasticities and implied coefficients. These issues complicate the production of consistent, commodity-level, provincial or national FAH coefficients suitable for food balance and policy analysis.

\subsection{The Gap Addressed by This Paper}

Existing studies have established that China's FAH consumption is substantial and growing and have identified correlates (income, demographics, urbanization) and data sources (urban household surveys, CHNS, expenditure surveys). However, they have not produced nationally or provincially representative, commodity-level out-of-home consumption \emph{coefficients} (quantity-based, total-to-at-home ratio) suitable for macro-level food demand and food security modeling. No existing study or international dataset provides such coefficients for China. This paper addresses the gap by: (1) using micro survey data (CHNS) to train machine-learning models for the at-home to total consumption ratio (and thus the out-of-home coefficient) by food category; (2) bridging micro and macro data via income-distribution matching so that predictions can be made for province-year cells with only macro covariates; (3) applying spatial interpolation and, where relevant, production structure to handle missing provinces; (4) comparing multiple ML models and reporting bootstrap-based prediction intervals; and (5) producing the first set of ML-based, commodity-level out-of-home consumption coefficients for China at provincial and national levels for use in food balance and food security assessments.
